\section{Conclusiones}

En este experimento se estudiaron circuitos RC y RLC mediante el análisis frecuencial de su respuesta, lo que permitió determinar la frecuencia de corte del primero y la frecuencia de resonancia del segundo. En el caso del circuito RC se estudió el comportamiento visto desde la resistencia (filtro PA) y el capacitor (filtro PB), en ambos se emplearon dos métodos independientes para estimar la frecuencia de corte \(f_0\). En primer lugar, se determinó la frecuencia de corte mediante la intersección de las asintotas del diagrama en escala logarítmica, correspondientes al régimen de funcionamiento del filtro y su zona de atenuación. Por otro lado, se ajustó el modelo teórico del desfasaje entre la señal de entrada y la respuesta en el resistor y el capacitor, considerando como parámetro libre la constante de tiempo $\tau$. Ambos resultados son compatibles entre sí, y a su vez, consistentes con el valor calculado con datos nominales \(f_{0_{ref}} = 1592(178)\) \SI{}{\hertz}, por lo que fueron promediados pesando su incerteza, obteniendo así \(f_0 = 1804(43)\) \SI{}{\hertz}. Esto respalda la validez del método de intersección de rectas para la caracterización de sistemas de primer orden.

Por otro lado, para el circuito RLC se determinó la frecuencia de resonancia \(f_r\) mediante el análisis del módulo de la función de transmisión en respuesta de la frecuencia visto desde la resistencia, como también mediante el desfasaje de la señal de la resistencia y la fuente. Los resultados son compatibles entre sí por lo que se realizó un promedio ponderado obteniendo \(f_r= 3318(19)\) \SI{}{\hertz}. A su vez, se estimó una resistencia efectiva adicional asociada a la presencia del inductor en el circuito de $R_L$ = \SI{157(9)}{\ohm}, así como una inductancia $L$ = 2,6(2) \(\times 10^{-2}\) \SI{}{\henry}. Lo que demuestra que el método empleado resulta adecuado no solo para determinar la frecuencia de resonancia, sino también para caracterizar parámetros internos del sistema. En última instancia, se graficaron los diagramas de Nyquist para ambos circuitos, los cuales resultaron compatibles con los gráficos usuales de un circuito RC y RLC.
